% Source: sources/returns/S6_04_MAPS_MANIFOLDS/sections/manifolds.tex:9012–9105
\paragraph{Theorem D.3.3 (Proved: the complex invariants of the smoothing).}
For $u\ne0$ sufficiently small and $0<|s|\ll1$, let $g:Y_{u,s}\to B$
be the above map and $S'_2=(g^{-1}(1))_{\mathrm{red}}$. Then
\[
 K_{Y_{u,s}}\simeq g^*\mathcal O_B(-1)\otimes\mathcal O_{Y_{u,s}}(2S'_2),
 \quad a(Y_{u,s})=1,
 \quad (h^{0,0},h^{0,1},h^{0,2},h^{0,3})=(1,1,0,0).
\]
In particular $K^{\otimes12}\simeq g^*\mathcal O_B(-6)$ and the
Kodaira dimension is $-\infty$.
\begin{proof}
We first compute the canonical bundle by actual transition forms.
On the core over a surface chart, the canonical form is
$\Omega_S\wedge da/a$. Changing a line frame adds a base one-form
to $da/a$, which wedges to zero with $\Omega_S$. On the second
plumbing collar the transition covers $T$ and changes the line
coordinate by the scalar $u\phi_s$; its logarithmic derivative
again contributes only a base form. Translation on the elliptic
surface preserves its invariant relative differential. Thus the
canonical transition functions are those of $K_S=\pi^*\mathcal O(-1)$.
On $xy=u(q-s)$ the residue expression
$dq\wedge dr\wedge dx/x$ becomes, using the equation,
a nonzero constant multiple of $dy\wedge dr\wedge dx$; it extends
across the node-chart loci where a line coordinate vanishes.
The expression over $q=zv$ has the same residue interpretation.
For $s\ne0$ it is regular wherever the total threefold is smooth.
Consequently moving $\phi$ introduces no new canonical divisor.

The two finite modifications give the same precise corrections as
in the original threefold. Indeed, for good reduction at $IV^*$
write $t=\zeta^3$, $x=\zeta^4X$, $y=\zeta^6Y$; then
$dt\wedge dx/y$ has order $2-2=0$ in $\zeta$. At $III$ write
$t-1=\zeta^4$, $x=\zeta^2X$, $y=\zeta^3Y$; its order is
$3-1=2$. Wedging with the logarithmic line coordinate gives the
threefold form, and the free quotient is etale as a map of total
spaces. These orders therefore descend respectively to zero and
$2S'_2$. This proves the asserted canonical formula. Since
$4S'_2=g^*(1)$, its twelfth power is $g^*\mathcal O(-6)$.
Connected fibres imply $g_*\mathcal O_Y=\mathcal O_B$, so none of
its positive powers has a section; if a power of $K$ had a section,
its twelfth power would contradict this fact.

For the algebraic dimension, use a marked regular fibre with
period matrix $\left(\begin{smallmatrix}6\mu&\tau&1&0\\
\beta_s&\mu&0&1\end{smallmatrix}\right)$.
Moving the zero changes only the scalar period $\beta_s$. Its
monodromy around the new zero $q=s$ is
$\beta_s\mapsto\beta_s\pm1$, while $\tau,\mu$ are unchanged.
If constants $a,b,c,d,e$ satisfy
$a+b\tau+c\mu+d\beta_s+e(6\mu^2-\tau\beta_s)=0$
identically, continuation around that zero gives $d-e\tau=0$.
The elliptic period $\tau$ is nonconstant, so $d=e=0$.
Continuation around infinity gives $b=0$ because
$\tau\mapsto\tau-1$ and $\mu$ is unchanged. Finally $\mu$ is
nonconstant: a constant satisfying
$\mu\circ g_1=(1-\mu)/\tau$ would force a nonconstant $\tau$ to
be constant, or give $0=1$. Hence $c=a=0$.
The type-$(1,1)$ equation for an integral alternating form now
shows that, outside a countable union of discrete subsets of the
base, the Neron--Severi group is generated by
$q=u\wedge w+6\gamma\wedge\delta$. Its associated Hermitian form
is indefinite: it is nondegenerate, its signature varies
continuously on the connected smooth-fibre locus, and at $s=0$
it has signature $(1,1)$ by the explicit period inequality.
It therefore supplies neither an ample class nor a nonzero
square-zero semipositive class. The algebraic reduction of a
complex torus is an abelian quotient; an abelian surface quotient
would supply the former class, and an elliptic quotient the latter.
Thus a very general fibre has algebraic dimension zero.
Every meromorphic function on $Y$ is constant on these fibres.
Its vertical differential consequently vanishes identically on
the regular locus, by the identity theorem, and it descends to a
meromorphic function on $B$. At an exceptional value this descent
extends meromorphically by restriction to a transverse holomorphic
disc and the bounded-pole criterion. Therefore
$\mathbb C(Y)=g^*\mathbb C(B)$ and $a(Y)=1$.

It remains to compute the structure-sheaf cohomology. The small
resolution $X^{\log}(u)\to Y_{u,0}$ has
$\pi_*\mathcal O_X=\mathcal O_{Y_{u,0}}$ and
$R^i\pi_*\mathcal O_X=0$ for $i>0$. For the latter assertion,
formal neighbourhoods of the exceptional curve have successive
quotients $\operatorname{Sym}^n(\mathcal O(1)^{\oplus2})$;
all their positive cohomology groups vanish, and formal functions
and coherence give the stated vanishing. Thus
$H^q(Y_{u,0},\mathcal O)=H^q(X^{\log}(u),\mathcal O)$, whose dimensions
are $(1,1,0,0)$, as computed by the supplied direct-image formulas.
Upper semicontinuity gives $h^{0,2}=h^{0,3}=0$ nearby. In a proper
flat analytic family the holomorphic Euler characteristic is
constant, as follows from the finite complex representing the
proper derived direct image over a small parameter disc. Its
value here is zero; connectedness gives $h^{0,0}=1$, and hence
$h^{0,1}=1$.
\end{proof}
